% !TeX program = pdflatex
% Valid document. Compile twice; see warnings-en.tex for deliberate errors.
\documentclass[11pt]{report}
\usepackage[T1]{fontenc}
\usepackage{amsmath}
\usepackage{xcolor}
\usepackage{listings}
\usepackage[hidelinks]{hyperref}
\setcounter{secnumdepth}{5}
\setcounter{tocdepth}{3}
\lstdefinestyle{sample}{
  basicstyle=\ttfamily\small,
  numbers=left,stepnumber=1,numberfirstline=true,numberblanklines=false,
  numbersep=8pt,frame=single,breaklines=true
}
\lstset{style=sample}
\title{LaTeX Structure Highlight: English Sample}
\author{Example document}
\date{}

\begin{document}
\maketitle
\tableofcontents
\chapter{A readable document}
\label{chap:overview}
\section{Experiment design}
\label{sec:design}
Each heading level has its own editor color. Compile twice to display actual
numbers above headings and references. Decorations do not change this PDF.

\subsection{Data preparation}
\label{sec:data}
Use \texttt{data\_file.csv} for a file name. The escaped underscore is valid.
The inline literal \verb|raw_name| is valid too.

\subsubsection{Input validation}
\label{sec:validation}
Math subscripts $x_i$ and \(y_j\) should not receive underscore warnings.
The sample enables numbering down to subparagraph with \texttt{secnumdepth=5}.

\paragraph{Implementation details}
\label{sec:implementation}
This numbered paragraph demonstrates another heading color.
\subparagraph{A smaller detail}
\label{sec:detail}
This heading is numbered too.
\subsection*{An unnumbered note}
Starred headings stay unnumbered and receive no compiled-number hint.

\section{Figures, tables and equations}
\label{sec:results}
See Figure~\ref{fig:flow}, Table~\ref{tab:counts}, and Equation~\eqref{eq:mean}.
An automatic reference is \autoref{fig:flow}. A citation is \cite{sample2026}.

\begin{figure}[htbp]
  \centering
  \fbox{\rule{0pt}{18mm}\rule{0.65\linewidth}{0pt}}
  \caption{A placeholder for the processing flow}
  \label{fig:flow}
\end{figure}
\begin{table}[htbp]
  \centering
  \caption{Example item counts}
  \label{tab:counts}
  \begin{tabular}{lr}
    Group & Count \\
    \hline
    Training & 80 \\
    Validation & 20 \\
  \end{tabular}
\end{table}
\begin{equation}
  \bar{x} = \frac{1}{n}\sum_{i=1}^{n}x_i
  \label{eq:mean}
\end{equation}
Paired display-math delimiters are valid:
\[
  \left(\sum_{i=1}^{n}x_i\right)^2 \geq 0.
\]

\section{Code and listing line numbers}
\label{sec:code}
The first block inherits the style and starts at ten. The language is explicit
on each block for grammar-based syntax colors.
\begin{lstlisting}[language=Python,firstnumber=10]
def normalize(value):
    cleaned = value.strip()

    return cleaned.lower()
\end{lstlisting}
Continue numbering, showing even numbers plus the first line:
\begin{lstlisting}[language=Python,firstnumber=last,stepnumber=2]
values = [" Alpha ", " BETA "]
normalized = [normalize(value) for value in values]
print(normalized)
\end{lstlisting}
Restart at one, with numbers on the right:
\begin{lstlisting}[language=Python,numbers=right,firstnumber=1]
for index, value in enumerate(normalized):
    print(index, value)
\end{lstlisting}

\begin{thebibliography}{9}
  \bibitem{sample2026}
  Example Author. \emph{A Fictional Guide to Readable Documents}. 2026.
  Sample data, not a real publication.
\end{thebibliography}

\begin{enumerate}
  \item First item.\label{item:first}
  \begin{enumerate}
    \item Nested item.
  \end{enumerate}
  \item Second item.
\end{enumerate}

\end{document}
